% ---------
% --------
%!TEX TS-program = pdflatex
%!TEX encoding = UTF-8 Unicode
%!TEX spellcheck = English (Aspell)

\documentclass[11pt]{article}
\usepackage{lucy,handout,graphicx}
\usepackage{fourier-orns}
\usepackage{mathtools}
\usepackage[normalem]{ulem} % either use this (simple) or
\usepackage{tikz}

\def\Cdot{\mathbin{\text{\normalfont \textbullet}}}
\def\Sym#1{\texttt{\color{BrickRed}#1}}
\def\BlueSym#1{\textbf{\texttt{\color{RoyalBlue}#1}}}
\allowdisplaybreaks

\begin{document}

\begin{center}
\Large\textbf{CSCI 332, Fall 2026}%
\\
\LARGE\textbf{Homework 3}%
\\[0.5ex]
\large Due Monday, September 14 Anywhere on Earth (6am Tuesday)
\end{center}

\bigskip
\hrule
\bigskip

\subsection*{Submission Requirements}
\begin{itemize}
    \item Type or clearly hand-write your solutions into a PDF format so that they are legible and professional. Submit your PDF on Gradescope. 
    \item Do not submit your first draft. Type or clearly re-write your solutions for your final submission. If your submission is not legible, we will ask you to resubmit.
    \item Use Gradescope to assign problems to the correct page(s) in your solution. If you do not do this correctly, we will ask you to resubmit.
\end{itemize}

\subsection*{Academic Integrity}

Remember, you may access \EMPH{any} resource in preparing your solution to the homework. However, you \EMPH{must}
\begin{itemize}
    \item write your solutions in your own words, and
    \item credit every resource you use (for example: ``Bob Smith helped me on
    this problem. He took this course at UM in Fall 2020''; ``I found a solution
    to a problem similar to this one in the lecture notes for a different
    course, found at this link: www.profzeno.com/agreatclass/lecture10''; ``Here is a link to the full transcript of my chat with ChatGPT about the problems on this HW''.)
    Ask if you need any clarifications about how to properly cite your sources!
\end{itemize}




\newpage
%----------------------------------------------------------------------

\headers{CSCI 332}{Homework 3 (due September 14)}{Fall 2026}


\begin{enumerate}
%\parindent 1.5em \itemsep 4ex plus 0.5fil

%----------------------------------------------------------------------

\item (0 points) List the resources you used for each problem in this homework. If you used no resources except
lecture notes and videos and the linked textbook, you can say ``none''. If you used
an AI tool, you should put a link to your full conversation here, and you must
use the
\href{https://umaal.org/files/teaching/csci332_fall2026/agents.md}{provided
prompt}
 to guide the AI to act
as a tutor, not an answer-provider.

You will be asked to resubmit your homework if this section is blank.

Remember that no matter what resources you consulted, you must write your solutions
\emph{in your own words}.

\item (3 points) Consider the number of the syllables in the pronunciation of the integer $n$.
For example, when $n=1$, the pronunciation is ``one'' and there is one syllable.
For each of the following, give the tightest/simplest function possible, or state that
there is no function possible.
\begin{enumerate}
    \item Give a function $f(n)$ so that the number of syllables in the pronunciation of $n$ is $O(f(n))$ for all $n$.
    \item Give a funciton $g(n)$ so that the number of syllables in the pronunciation is of $n$ is $\Omega(g(n))$ for all $n$.
    \item Give a function $h(n)$ so that the number of syllables in the pronunciation is of $n$ is $\Theta(h(n))$ for all $n$.
\end{enumerate}

\item (3 points) Prove that $\log_a n = \Theta(\log_b n)$ for any positive constants $a, b
> 1$ by proving each of the following. You must use the definition of big-O and
big-Omega in your proofs; that is, you must give a $c$ and an $n_0$. (Hint: Use the change of base formula for logarithms,
which states that $\log_a n = \frac{\log_b n}{\log_b a}$.)

    \begin{enumerate}[(a)]
        \item $log_a n = O(log_b n)$
        \item $log_a n = \Omega(log_b n)$
    \end{enumerate} 

\item (4 points) For each of the following algorithms, you will analyze both the worst and best case runtimes.
\begin{enumerate}
    \item
    \begin{algo}
        \underline{IsPrime(n):}\\
        \> if $n < 2$\\
        \> \> return False\\
        \> for $d = 2$ to $\lfloor \sqrt{n} \rfloor$\\
        \> \> if $n \bmod d = 0$\\
        \> \> \> return False\\
        \> return True
    \end{algo}
    \begin{enumerate}
        \item Describe a best-case input to IsPrime.
        \item Give a function $f(n)$ so that the best-case runtime of IsPrime is $\Theta(f(n))$.
        \item Describe a worst-case input to IsPrime.
        \item Give a function $f(n)$ so that the worst-case runtime of IsPrime is $\Theta(f(n))$.
    \end{enumerate}
    \item 
    \begin{algo}
        \underline{HasDuplicate(A[1..n]):}\\
        \> for $i = 1$ to $n$\\
        \> \> for $j = i+1$ to $n$\\
        \> \> \> if $A[i] = A[j]$\\
        \> \> \> \> return True\\
        \> return False
    \end{algo}
    \begin{enumerate}
        \item Describe a best-case input to HasDuplicate.
        \item Give a function $f(n)$ so that the best-case runtime of HasDuplicate is $\Theta(f(n))$.
        \item Describe a worst-case input to HasDuplicate.
        \item Give a function $f(n)$ so that the worst-case runtime of HasDuplicate is $\Theta(f(n))$.
    \end{enumerate}
\end{enumerate}



\end{enumerate}
%%%%%%%%%%%%%%%%%%%%
\end{document}
